\begin{helper}
Let \(V\) be finite, let \(u\ne v\), and let
\(B_U,B_T\subseteq V\setminus\{u,v\}\) be disjoint.  If
\(0\le b\le a\le1\), then in the complementary coordinate space
\(\mathbb R^{V\setminus\{u,v\}}\) the product box with side interval
\([0,a)\) on \(B_U\), side interval \([0,b)\) on \(B_T\), and side interval
\([0,1)\) on all other complementary coordinates has volume
\[
a^{|B_U|}b^{|B_T|}.
\]
\end{helper}
\begin{proof}
Let \(W=V\setminus\{u,v\}\).  Define the upper endpoint
\[
U(w)=
\begin{cases}
a,&w\in B_U,\\
b,&w\in B_T,\\
1,&w\notin B_U\cup B_T.
\end{cases}
\]
The box in the statement is \(\prod_{w\in W}[0,U(w))\).  By the finite
product formula for half-open intervals, its Lebesgue volume is
\[
\prod_{w\in W} (U(w)-0).
\]
Since \(B_U\) and \(B_T\) are disjoint and both avoid \(u\) and \(v\), their
images inside \(W\) have exactly the same cardinalities as \(B_U\) and
\(B_T\).  The product over \(W\) therefore has one factor \(a\) for each
element of \(B_U\), one factor \(b\) for each element of \(B_T\), and one
factor \(1\) for every remaining coordinate.  Hence the product is
\(a^{|B_U|}b^{|B_T|}\), as claimed.
\end{proof}
